\section*{Dimension-free alpha-log payment on adaptive restriction trees}

All logarithms are natural.  Let \(h:\{-1,1\}^n\to\{-1,1\}\) be an
arbitrary Boolean sign.  Consider any deterministic restriction tree.
At an internal node \(v\), the path has fixed a partial assignment
\(\rho_v\), the node queries a coordinate \(i(v)\notin\operatorname{dom}
ho_v\),
and no root-to-leaf path queries a physical coordinate twice.  Write
\[
 w_v=2^{-|\operatorname{dom}\rho_v|},\qquad
 h_v=h\vert_{\rho_v},\qquad
 \alpha_v=\widehat h_v(\{i(v)\}).
\]

\subsection*{Theorem (proved): adaptive nonreusing hierarchy payment}

For every such Boolean sign and every such tree,
\[
 \boxed{
 \sum_{v\ \mathrm{internal}}w_v\alpha_v^2
       \log\frac1{\alpha_v^2}
 \le\frac2e\,\mathbf I(h).}                                      \tag{1}
\]
In bits the coefficient is \(2/(e\log2)<1.062\).

\paragraph{Proof.}
At every node, the Boolean edge argument gives
\[
 |\alpha_v|\le\operatorname{Inf}_{i(v)}(h_v).
\]
Since \(u^2\log(1/u^2)\le(2/e)u\) on \([0,1]\),
\[
 w_v\alpha_v^2\log\frac1{\alpha_v^2}
 \le\frac2e w_v\operatorname{Inf}_{i(v)}(h_v).                    \tag{2}
\]
Fix a physical coordinate \(i\).  The cylinder events corresponding to
nodes which query \(i\) are pairwise disjoint, because no input path can
query \(i\) twice.  Moreover, the path preceding such a node has not
looked at \(x_i\).  Flipping \(x_i\) therefore preserves the event of
reaching that node, and
\[
 w_v\operatorname{Inf}_i(h_v)
 =\Pr\{X\text{ reaches }v\text{ and }i\text{ is pivotal for }h\}.
\]
Summing over the disjoint nodes which query \(i\) gives at most
\(\operatorname{Inf}_i(h)\).  Sum this inequality over \(i\), then use
(2), proving (1). \(\square\)

\subsection*{Composition corollary and exact scope}

If an outer Boolean pair \(a\pm b\) remains fixed while a physical
restriction tree is run inside a disjoint child block, put
\(q=\|b\|_2^2\).  Multiplying (1) by \(q\) yields
\[
 \sum_v q w_v\alpha_v^2\log\frac1{\alpha_v^2}
 \le\frac2e q\mathbf I(h).                                       \tag{3}
\]
The right side is the exact influence contribution of that child block
at the balanced composition interface.  Thus rare restriction nodes,
adaptive coordinate choices, and arbitrarily large tree depth do not
accumulate a depth factor for the universal alpha-log envelope.

Equation (3) does not by itself cover a hierarchy in which the outer
activation weight changes coherently with the same restricted physical
variables.  Nor does it upper-bound the complete biased remainder
\(R_\mu(\alpha)\).  Those are the remaining joint-witness and coherent
interference obligations.

\subsection*{Mandatory controls}

The theorem is simultaneous over all nodes, so it survives:
\begin{itemize}
 \item parity-zero erasure, since \(\alpha_v=0\) contributes zero;
 \item rare restrictions, because each term carries its exact cylinder
       probability \(w_v\);
 \item adaptive addressing trees, because the queried coordinate may
       depend on all earlier answers;
 \item near tensors and prime Boolean gates, since no decomposition of
       \(h\) is assumed;
 \item arbitrary depth, because pivotal events are grouped by physical
       coordinate rather than paid once per level.
\end{itemize}
It makes no overlapping-input composition claim.  Shared-variable
interference remains governed by the separate overlap ledger.

